\section{Các công trình liên quan}

\subsection{Truy xuất thông tin đặc và Đa vector}
Truy xuất đặc sử dụng kiến trúc bộ mã hóa kép (Bi-Encoder) đã thay đổi mạnh mẽ lĩnh vực truy xuất thông tin hiện đại~\cite{reimers2019sentencebert,karpukhin2020dense}. Các mô hình nhúng đa ngôn ngữ ngữ cảnh dài như BGE-M3~\cite{chen2024bge} hỗ trợ cửa sổ lên tới $8{,}192$ token và tín hiệu đa độ mịn. Tuy nhiên, với tài liệu dài có cấu trúc bán phần như CV, việc nén các phân mục không đồng nhất vào một vector duy nhất vẫn gây loãng ngữ nghĩa, còn các mô hình nhẹ chạy tốt trên CPU như \texttt{all-MiniLM-L6-v2} bị cắt cụt sau $256$ token. Các kiến trúc tương tác muộn cấp token như ColBERT~\cite{khattab2020colbert} giữ được độ mịn qua phép toán MaxSim nhưng làm phình chỉ mục $30\times$--$100\times$ và tăng mạnh chi phí ma trận trên CPU. Gần đây, \v{C}eovi\'{c} và cộng sự~\cite{ceovic2026multiview} đề xuất ghép cặp ứng viên--công việc đa góc nhìn dùng học tương phản và tăng cường diễn giải bằng LLM. Khác với các hướng đó, phương pháp nhận biết phân mục của chúng tôi vận hành ở mức phân mục chức năng thô ($M \le 4$ vector/tài liệu), tránh phình chỉ mục cấp token và chi phí sinh dữ liệu LLM, đồng thời chống cắt cụt và hỗ trợ giải trình theo từng phân mục.

\subsection{Ghép cặp Ứng viên--Công việc và Hệ phân loại Kỹ năng}
Mô hình hóa mức độ phù hợp nhân sự--công việc đã tiến hóa từ truy xuất túi từ sang học biểu diễn sâu và mạng nơ-ron đồ thị~\cite{zhu2018person,bian2020learning,zhang2021personalized}. Các đánh giá công nghiệp cho thấy thế mạnh bổ trợ giữa vector đặc và BM25~\cite{robertson2009probabilistic} khi dung hợp qua Reciprocal Rank Fusion (RRF)~\cite{cormack2009reciprocal}, như khảo sát gần đây về ghép cặp ứng viên ngữ nghĩa doanh nghiệp của Yashwant và cộng sự~\cite{yashwant2026semantic}. Tuy nhiên, so khớp tập kỹ năng truyền thống thường dùng Jaccard đối xứng, vô tình phạt ứng viên giàu kinh nghiệm sở hữu nhiều kỹ năng vượt ngoài yêu cầu trực tiếp của JD. Công trình của chúng tôi giải quyết vấn đề này bằng độ phủ yêu cầu bất đối xứng kết hợp đồ thị kỹ năng có kiểu quan hệ đa bước (phân biệt quan hệ bao hàm, framework và hệ sinh thái đồng hành, đồng thời chặn tuyệt đối cộng điểm sai giữa hai kỹ năng không liên kết).

\subsection{Xếp hạng lại Đa giai đoạn và Học xếp hạng Đơn điệu}
Các pipeline truy xuất độ chính xác cao thường áp dụng thiết kế hai giai đoạn nối tầng~\cite{asadi2013effectiveness,nogueira2019passage}, kết hợp truy hồi nhanh ở tầng 1 với bộ xếp hạng lại Cross-Encoder ở tầng 2. Các hệ thống gần đây như ConFit~v3~\cite{jasonyux2026confitv3} dùng học tăng cường và chưng cất danh sách trên LLM mã nguồn mở (Qwen3-8B/32B) để xếp hạng lại CV dài. Dù hiệu quả trên cụm GPU, Cross-Encoder và bộ xếp hạng lại LLM chịu độ trễ rất lớn trên CPU và dễ suy giảm do dịch chuyển miền khi huấn luyện trên dữ liệu hỏi đáp đoạn ngắn như MS~MARCO. Hơn nữa, các bộ xếp hạng nơ-ron không ràng buộc thiếu bảo đảm đơn điệu cấu trúc. Chúng tôi kiểm chứng trực tiếp đánh đổi tính toán--độ chính xác này trước Calibrated Hybrid một giai đoạn, đồng thời đánh giá Học xếp hạng ràng buộc đơn điệu (Linear LTR chiếu đơn hình và LambdaMART đơn điệu) trên tập đặc trưng phân mục, từ khóa và taxonomy.
